\scanpage{110}
IV 7 Oubli au n\textsuperscript{o} 7.6.

\uline{Proposition} 6.7. Soient $f:X\longrightarrow Y$ un morphisme
\deleted{de présentation finie} \deleted{et} propre,
\inserted{avec $Y$ loc. noeth.,}
$F$ un Module \deleted{de présentation finie} \inserted{cohérent}
sur $X$. \inserted{de support $X$,} Pour tout $y\quad Y$,
Soit $\ell(y)$ la somme des \deleted{long} multiplicités \deleted{tota}
géométriques totales \inserted{sur $k(y)$} de
$F_y$ en les points maximaux de $\operatorname{supp}F_y$.
\inserted{On suppose $f$ univ. ouvert en les pts maximaux de $X_{y_0}$,
et $F_{y_0}$ sans \unsure{idéaux} premiers ass. immergés}
Alors la fonction $\unsure{z|y}\mapsto\ell(\unsure{z|y})$ sur
$Y$ est semi-continue supérieurement\inserted{en $y_0$}.
\margin{?}

Le procédé breveté nous ramène au cas où $Y$ est noethérien. On sait
(réf) que la fonction $\ell$ est constructible, il reste donc (réf) à prouver
que $y\in\overline{y'}$ implique $\ell(y)\geq\ell(y')$. Le procédé habituel nous permet
alors de nous ramener au cas où $Y$ est le spectre d'un anneau de valuation
discrète, dont $y$ ($y'$) est le point fermé (générique), et \deleted{\unsure{à prouver}}
à prouver dans ce cas l'inégalité correspondante sur les longueurs:

\uline{Lemme} 6.8. Soit $f:X\quad Y$ un morphisme propre, avec $Y$ spectre d'un anneau
de valuation discrète, $F$ un Module cohérent sur $X$, \inserted{de support $X$,}
$F_0$ et $F_1$ les \deleted{fibr}
Modules induits sur la fibre spéciale resp. la fibre générique, $x_i$
($x'_j$) les points maximaux de $\operatorname{supp}F_0$ ($\operatorname{supp}F_1$).
\inserted{(Supposons que tt comp. de $X$ domine $Y$ et $F_0$ soit $(S_1)$.)}
Alors on a
\[
\operatorname{long} F_{0\,x_i}\geq\sum_i\operatorname{long}F_{1\,x'_j}
\]
Si l'inégalité est vérifiée, alors $F$ est plat sur $Y$ aux points \deleted{\unsure{génériques}}
maximaux de $\operatorname{supp}F_0$, \deleted{donc toute composante de $\operatorname{supp}F$ domine $Y$,}

Soit $t$ une uniformisante pour $V=\underline O_{Y,y}$, soit $N$ le sous-Module
de $F$ réunion des annulateurs des $t^n$, donc on a une suite exacte
\[
0\mathrel{-}N\mathrel{-}F\mathrel{-}G\mathrel{-}0
\]
d'où une suite exacte (grâce au fait que $t$ est $G$-régulier):
\[
0\mathrel{-}N/tN\mathrel{-}F/tF\mathrel{-}G/tG\mathrel{-}0
\]
\unsure{Qela|Cela} montre que les deux sommes \deleted{pour} relatives à $F$ qu'il s'agit de
comparer, s'obtiennent en ajoutant les sommes analogues pour $N,G$.
Comme $N_1=0$, pour prouver notre inégalité, il suffit donc de
prouver l'inégalité analggue pour $G$; et on ne peut avoir égalité
que si $N/tN$ \deleted{est} $=N_0$ est nul aux points maximaux de $\operatorname{supp}$
\unsure{$\cdots$}

\margin{\figure{\begin{tikzpicture}[x=.5cm,y=.5cm]
\draw (0,0)..controls(-.1,2)and(.3,9)..(.9,11);
\draw (4,-1)..controls(4.7,1)and(3.5,8)..(6,15);
\draw (8,0)..controls(8.5,3)and(6.8,7)..(7.5,12);
\end{tikzpicture}}}

\margin{N.B. J'ignore si les assertions sont vraies si on ne
suppose pas déjà $F$ plat$/Y$ aux pts maximaux
de $X_y$. Question liée: $A$ anneau local noeth., $M$ module de
type fini sur $A$, $t\in A$, $\unsure{t\notin\mathfrak p}$\deleted{\unsure{minimal}} pour tt $\mathfrak p$ associé,
on suppose $\operatorname{supp}M=\operatorname{Spec}A$, $\dim\unsure{A/\mathfrak p_0}=1$,
prouver $\mathfrak m\in\operatorname{Ass}M/tM$ !}
\mspage{109}

